\documentclass[8pt,landscape,a4paper]{article}
\usepackage[utf8]{inputenc}
\usepackage[landscape,margin=0.4in]{geometry}
\usepackage{multicol}
\usepackage{xcolor}
\usepackage{listings}
\usepackage{amsmath}
\usepackage{titlesec}
\usepackage{enumitem}
\usepackage{fancyhdr}

% Reduce spacing
\setlength{\parindent}{0pt}
\setlength{\parskip}{1pt}
\setlist{nosep, leftmargin=*, itemsep=0pt}

% Section formatting
\titleformat{\section}{\normalsize\bfseries\color{blue!70!black}}{\thesection}{0.4em}{}
\titlespacing*{\section}{0pt}{3pt}{1pt}
\titleformat{\subsection}{\small\bfseries\color{blue!50!black}}{\thesubsection}{0.4em}{}
\titlespacing*{\subsection}{0pt}{2pt}{0.5pt}

% Code style
\lstset{
    language=Python,
    basicstyle=\ttfamily\tiny,
    keywordstyle=\color{blue!70!black},
    commentstyle=\color{green!50!black},
    stringstyle=\color{red!70!black},
    showstringspaces=false,
    breaklines=true,
    frame=single,
    rulecolor=\color{gray!50},
    xleftmargin=1pt,
    xrightmargin=1pt,
    aboveskip=1pt,
    belowskip=1pt,
    columns=flexible
}

\pagestyle{fancy}
\fancyhf{}
\renewcommand{\headrulewidth}{0.4pt}
\fancyhead[L]{\textbf{CMPUT-274 Python \& Functional Programming Cheat Sheet}}
\fancyhead[R]{Page \thepage}

\begin{document}

\begin{multicols}{3}
\setlength{\columnseprule}{0.4pt}

\section*{Python Basics}

\subsection*{Data Types \& Operations}
\begin{lstlisting}
# Numbers
x = 42          # int
y = 3.14        # float
z = x + y       # arithmetic
w = x // 10     # integer division
exp = 2**3      # exponentiation (8)

# Strings
s = "hello"
s[0]            # 'h' (indexing)
s[1:]           # 'ello' (slicing)
s + "world"     # concatenation
len(s)          # 5
\end{lstlisting}

\subsection*{Control Structures}
\begin{lstlisting}
# If statements
if x < 0:
    return "negative"
elif x == 0:
    return "zero"
else:
    return "positive"

# Assertions
assert n >= 1
\end{lstlisting}

\section*{Recursion Basics}

\subsection*{Simple Recursion Pattern}
\textbf{Base case} + \textbf{Recursive case}

\begin{lstlisting}
def factorial(n):
    if n <= 1:      # Base case
        return 1
    # Recursive case
    return n * factorial(n-1)

# factorial(5) -> 120
\end{lstlisting}

\subsection*{String Recursion}
\begin{lstlisting}
def strlen(s):
    if s == "":     # Base case
        return 0
    return 1 + strlen(s[1:])

# strlen("hello") -> 5
\end{lstlisting}

\subsection*{Numeric Recursion}
\begin{lstlisting}
def logFloor(n):
    # floor(log_10(n))
    assert n >= 1
    if n < 10:
        return 0
    return 1 + logFloor(n/10)

def digitSum(n):
    if n < 10:
        return n
    digit = n//(10**logFloor(n))
    remainder = n - digit*10**logFloor(n)
    return digit + digitSum(remainder)

# digitSum(456) -> 15
\end{lstlisting}

\section*{Linked Lists (LList)}

\subsection*{LList Operations}
\begin{lstlisting}
from cmput274 import *

# Create empty list
l = empty()

# Construct list (prepend)
l = cons(1, cons(2, cons(3, empty())))
# or use shorthand:
l = LL(1, 2, 3)  # -> (1, 2, 3)

# Access operations
first(l)         # 1 (head)
rest(l)          # (2, 3) (tail)
isEmpty(l)       # False
\end{lstlisting}

\subsection*{LList Recursion}
\begin{lstlisting}
def doubleList(l):
    if isEmpty(l):
        return l
    return cons(2*first(l), 
                doubleList(rest(l)))

# doubleList(LL(1,2,3)) -> (2,4,6)
\end{lstlisting}

\subsection*{List Building}
\begin{lstlisting}
def ascendList(n):
    # Build [0,1,2,...,n]
    if n == 0:
        return cons(0, empty())
    return append(n, ascendList(n-1))

def append(elem, l):
    if isEmpty(l):
        return cons(elem, l)
    return cons(first(l), 
                append(elem, rest(l)))
\end{lstlisting}

\subsection*{List Access}
\begin{lstlisting}
def getIth(l, i):
    # Get ith element (0-indexed)
    if i == 0:
        return first(l)
    return getIth(rest(l), i-1)

# myL = LL(5, 10, 27)
# getIth(myL, 1) -> 10
\end{lstlisting}

\columnbreak

\section*{Higher-Order Functions}

\subsection*{Map (Transform)}
\begin{lstlisting}
def map(f, l):
    # Apply f to each element
    if isEmpty(l):
        return empty()
    return cons(f(first(l)), 
                map(f, rest(l)))

# map(lambda x: x+3, LL(1,2,3))
# -> (4, 5, 6)
\end{lstlisting}

\subsection*{Filter (Select)}
\begin{lstlisting}
def filter(p, l):
    # Keep elements where p(e) is True
    if isEmpty(l):
        return empty()
    ror = filter(p, rest(l))
    if p(first(l)):
        return cons(first(l), ror)
    return ror

# def isNeg(x): return x < 0
# filter(isNeg, LL(-10,5,-2,0))
# -> (-10, -2)
\end{lstlisting}

\subsection*{Fold (Reduce)}
\begin{lstlisting}
def fold(l, combine, base):
    # Right fold
    if isEmpty(l):
        return base
    return combine(first(l), 
        fold(rest(l), combine, base))

# Sum example
def add(x, y): return x + y
# fold(LL(1,2,3), add, 0) -> 6

# Count ones in binary string
def count1(c, n):
    if c == "1":
        return 1 + n
    return n
# foldr("1011", count1, 0) -> 3
\end{lstlisting}

\subsection*{Build List with Function}
\begin{lstlisting}
def buildList(f, n):
    # Apply f to [0..n-1]
    return map(f, ascendList(n-1))

# def pow2(x): return 2**x
# buildList(pow2, 5)
# -> (1, 2, 4, 8, 16)
\end{lstlisting}

\section*{Accumulator Pattern}

\subsection*{Tail Recursion (Optimized)}
\begin{lstlisting}
@TCO
def ascendListAcc(n, acc):
    if n == 0:
        return cons(0, acc)
    return ascendListAcc(n-1, 
                         cons(n, acc))

def ascendList(n):
    return trampoline(
        ascendListAcc(n, empty()))
\end{lstlisting}

\textbf{Benefits:} Space-efficient, avoids stack overflow for large inputs.

\subsection*{Finding Extremes}
\begin{lstlisting}
def exComb(n, pair):
    # Update (MIN, MAX)
    if n < first(pair):
        return cons(n, rest(pair))
    if n > first(rest(pair)):
        return LL(first(pair), n)
    return pair

def extremes(l):
    return foldl(l, exComb, 
        LL(first(l), first(l)))

# extremes(LL(5,2,9,1)) -> (1,9)
\end{lstlisting}

\section*{String Manipulation}

\subsection*{Pig Latin}
\begin{lstlisting}
def pl(s):
    vowels = "aeiou"
    if s[0] in vowels:
        return s + "way"
    else:
        return s[1:] + s[0] + "ay"

# pl("hello") -> "ellohay"
# pl("orange") -> "orangeway"
\end{lstlisting}

\columnbreak

\subsection*{Leet Speak}
\begin{lstlisting}
def convertChar(c):
    conversions = {
        'a':'4', 'A':'4',
        'l':'1', 'L':'1',
        't':'7', 'T':'7',
        'e':'3', 'E':'3',
        's':'5', 'S':'5',
        'o':'0', 'O':'0'
    }
    return conversions.get(c, c)

def leetSpeak(s):
    if s == "":
        return ""
    if s[:2].lower() == "er":
        return "z0r" + leetSpeak(s[2:])
    return convertChar(s[0]) + \
           leetSpeak(s[1:])
\end{lstlisting}

\subsection*{String Equality}
\begin{lstlisting}
def streq(s1, s2):
    if s1 == "" and s2 == "":
        return True
    if s1 == "" or s2 == "":
        return False
    if s1[0] != s2[0]:
        return False
    return streq(s1[1:], s2[1:])
\end{lstlisting}

\section*{Common Patterns}

\subsection*{List Operations}
\begin{lstlisting}
def sum(lon):
    if isEmpty(lon):
        return 0
    return first(lon) + sum(rest(lon))

def length(lon):
    if isEmpty(lon):
        return 0
    return 1 + length(rest(lon))

def mean(lon):
    return sum(lon) / length(lon)
\end{lstlisting}

\subsection*{Parity Check}
\begin{lstlisting}
def parity(s):
    # True if even number of 1s
    def count1(c, n):
        return (1+n) if c=="1" else n
    return foldr(s, count1, 0) % 2 == 0

# Alternative with boolean
def checkOne(c, b):
    return not b if c=="1" else b
def parity2(s):
    return foldr(s, checkOne, True)
\end{lstlisting}

\section*{Efficiency Tips}

\subsection*{Time Complexity}
\begin{itemize}[itemsep=0pt]
    \item \textbf{Linear recursion:} $O(n)$ time
    \item \textbf{Tree recursion:} Can be $O(2^n)$ 
    \item \textbf{List traversal:} $O(n)$
    \item \textbf{Append to end:} $O(n)$ (prepend is $O(1)$)
\end{itemize}

\subsection*{Space Complexity}
\begin{itemize}[itemsep=0pt]
    \item \textbf{Regular recursion:} $O(n)$ stack space
    \item \textbf{Tail recursion + TCO:} $O(1)$ stack
    \item \textbf{Accumulator pattern:} More efficient
\end{itemize}

\subsection*{Fast vs Slow}
\begin{lstlisting}
# SLOW: Append to end each time
def ascendListSlow(n):
    if n == 0:
        return cons(0, empty())
    return append(n, ascendListSlow(n-1))
# O(n^2) - appending is O(n)

# FAST: Use accumulator
def ascendListFast(n, acc=empty()):
    if n == 0:
        return cons(0, acc)
    return ascendListFast(n-1, cons(n, acc))
# O(n) - prepending is O(1)
\end{lstlisting}

\section*{Key Concepts}

\subsection*{Functional Programming}
\begin{itemize}[itemsep=0pt]
    \item \textbf{Immutability:} Don't modify data
    \item \textbf{First-class functions:} Pass functions as arguments
    \item \textbf{Pure functions:} No side effects
    \item \textbf{Recursion:} Instead of loops
\end{itemize}

\subsection*{Design Recipe}
\begin{enumerate}[itemsep=0pt]
    \item \textbf{Signature:} Input/output types
    \item \textbf{Purpose:} What it does
    \item \textbf{Examples:} Test cases
    \item \textbf{Body:} Implementation
    \item \textbf{Test:} Verify examples
\end{enumerate}

\subsection*{Recursion Template}
\begin{lstlisting}
def recursive_fn(data):
    # 1. Base case(s)
    if base_condition:
        return base_value
    
    # 2. Recursive case
    # - Make problem smaller
    # - Recurse on smaller problem
    ror = recursive_fn(smaller_data)
    
    # 3. Combine results
    return combine(current, ror)
\end{lstlisting}

\end{multicols}

\end{document}
